\chapter{Operating a CL Development Environment}

This chapter covers the hands-on side of working with Common Lisp:
installing it, talking to it, connecting it to your editor, loading
libraries, and finding out what went wrong when something does. If
Lisp code is completely new to you, you may want to glance ahead at
Chapter~3 as you read.

It helps to think of a running Common Lisp as a small operating system
of its own, sitting on top of whichever operating system your machine
runs. You start it, and then you load things into it, ask it
questions, and change it while it runs. People once built whole
workstations this way, with Lisp all the way down.

Have a Lisp running while you read, and type the examples in.

\section{Choosing an Implementation}

Common Lisp is a standard with several implementations\index{implementations}, in the way
that C and SQL are. The ones you are most likely to meet:

\begin{itemize}
\item \textbf{SBCL}\index{SBCL} (Steel Bank Common Lisp), \url{https://www.sbcl.org}.
Free and open source, with a compiler that produces fast native code
and gives thorough warnings. The most widely used implementation.

\item \textbf{CCL}\index{CCL} (Clozure CL), \url{https://ccl.clozure.com}. Free and
open source. Compiles very quickly, and is pleasant to develop in.

\item \textbf{ECL}\index{ECL} (Embeddable Common Lisp),
\url{https://ecl.common-lisp.dev}. Free. Compiles by way of C, so
it can be linked into C and C++ programs and runs almost anywhere.

\item \textbf{ABCL}\index{ABCL} (Armed Bear Common Lisp), \url{https://abcl.org}.
Free. Runs on the Java virtual machine and can call Java libraries.

\item \textbf{Allegro CL}\index{Allegro CL}, from Franz Inc., \url{https://franz.com}, and
\textbf{LispWorks}\index{LispWorks}, \url{https://www.lispworks.com}. Commercial
products with their own development environments, graphical toolkits
and support, each with a free edition for learning.
\end{itemize}

They all implement the same language, and what you learn on one
carries over to the others. This book uses SBCL. Where something in
the text is specific to SBCL, it says so.

\section{Installing a CL Environment}

SBCL is packaged for the common systems:

\begin{verbatim}
sudo apt install sbcl         (Debian, Ubuntu)
brew install sbcl             (macOS with Homebrew)
\end{verbatim}

\noindent
For Windows there is an installer at \url{https://www.sbcl.org}. The
packaged versions can be a few releases behind; the website has
current binaries if you want them.

If you use Docker\index{Docker}, you can try an implementation without installing
anything. The \code{clfoundation} images on Docker Hub cover most of
the free implementations:

\begin{verbatim}
docker run -it --rm clfoundation/sbcl
\end{verbatim}

\section{Running CL in a Terminal}

The simplest way to run CL is from a terminal. Few people work this
way all day, because an editor connected to Lisp is far more
convenient, but it is the place to start, and it is still useful for
logging in to a Lisp that is running on a remote machine.

\subsection{Starting and Stopping}

To start SBCL, type

\begin{verbatim}
sbcl
\end{verbatim}

\noindent
You will see a few lines of introduction and then a prompt:

\begin{verbatim}
This is SBCL 2.6.9, an implementation of ANSI Common Lisp.
More information about SBCL is available at <http://www.sbcl.org/>.
...
*
\end{verbatim}

CL has entered its read-eval-print loop\index{REPL} and is waiting for you to type
something. It will read what you type, evaluate it to get a return
value, and print that value:

\begin{verbatim}
* (+ 2 2)
4
\end{verbatim}

To leave SBCL, type

\begin{verbatim}
(exit)
\end{verbatim}

\noindent
Each implementation has its own name for this (in CCL it is
\code{(quit)}). Typing the end-of-file character, Control-D, works in
most of them.

The bare terminal prompt has no command history and no line editing.
If you plan to spend any time there, install \code{rlwrap} and start
Lisp as \code{rlwrap sbcl}.

\subsection{Compiling and Loading a File}

\textbf{Try it:} With your text editor, create a file
\code{/tmp/hello.lisp} containing the following (do not worry about
understanding it yet):

\begin{verbatim}
(defun hello ()
  (write-string "Hello, World!"))
\end{verbatim}

Now, at the CL prompt, compile and load the file. The text after each
\code{*} prompt is what you type; everything else is printed by CL:

\begin{verbatim}
* (compile-file "/tmp/hello.lisp")
; compiling file "/tmp/hello.lisp" (written 05 OCT 2026 06:48:14 PM):

; wrote /tmp/hello.fasl
; compilation finished in 0:00:00.000
#P"/tmp/hello.fasl"
NIL
NIL
* (load "/tmp/hello")
T
\end{verbatim}

\code{Compile-file}\index{compile-file@\texttt{compile-file}} translates the source file into a binary file of
machine code. The binary is called a \emph{fasl} file\index{fasl file}, for ``fast
loading.'' \code{Load}\index{load@\texttt{load}} brings a file's definitions into the running
Lisp. Given a name with no extension, it picks the fasl file if there
is one.

You can also \code{load} a source file directly, and the effect is the
same, except that it has to be processed each time. Fasl files are
specific to the implementation, its version, and the kind of machine.
You do not pass them around; each Lisp makes its own from the source.

Now call the function you defined:

\begin{verbatim}
* (hello)
Hello, World!
"Hello, World!"
\end{verbatim}

The text appears twice. The first line is the function printing the
string, as you asked it to. The second is the REPL printing the value
that the function returned, which happens to be the same string; it
is shown in double quotes because the REPL prints values in the form
you would type them. Chapter~3 has more on this.

The point for now is the way of working: definitions live in files,
you compile and load them into a running Lisp, and you call them from
the prompt.

\section{Running CL inside an Editor}

The way most people develop in CL is with the Lisp running as a
subprocess of their text editor. The editor then knows about the
running Lisp. With a keystroke you can send it the function you are
editing, ask it for a function's arguments or documentation, jump to
the source of anything that is defined, and look inside live data.

\subsection{Emacs with SLIME or Sly}

The traditional editor for Lisp is GNU Emacs\index{Emacs}, which is itself largely
written in a dialect of Lisp, and the standard connection between
Emacs and Common Lisp is SLIME\index{SLIME} (\url{https://slime.common-lisp.dev}).
Sly\index{Sly} is a fork of SLIME with some added conveniences. Both work with
every implementation listed above, and what follows applies to both.

If you are new to Emacs, run its built-in tutorial first (it is under
the Help menu). Emacs commands are written as key chords: \code{C-x}
means hold Control and press x; \code{M-x} means hold Meta, which is
the Alt or Option key, and press x.

To install SLIME, add the MELPA package archive to Emacs if you have
not already (\url{https://melpa.org} has the two lines to copy), then
type \code{M-x package-install}, press Return, and give the name
\code{slime}. Then put this line in your Emacs init file to say
which Lisp to run:

\begin{verbatim}
(setq inferior-lisp-program "sbcl")
\end{verbatim}

Now \code{M-x slime} starts Lisp and opens a REPL in an Emacs window.
The prompt shows the current package, a notion explained in
Section~\ref{sec:packages}:

\begin{verbatim}
CL-USER> (+ 2 2)
4
\end{verbatim}

\noindent
The rest of this book shows examples with this prompt.

Table~\ref{tab:slime} lists the commands you will use most.

\begin{table}[ht]
\centering
\begin{tabular}{ll}
\hline
\textbf{Key} & \textbf{Action} \\
\hline
\code{C-c C-c} & Compile the definition the cursor is in \\
\code{C-c C-k} & Compile and load the whole file \\
\code{C-x C-e} & Evaluate the expression before the cursor \\
\code{C-c C-z} & Switch to the REPL \\
\code{M-.}     & Go to the definition of the name at the cursor \\
\code{M-,}     & Come back from there \\
\code{C-c C-d d} & Describe the name at the cursor \\
\code{C-c RET} & Expand the macro call at the cursor \\
\code{M-p}, \code{M-n} & In the REPL: previous and next input \\
\code{C-c C-c} & In the REPL: interrupt a running computation \\
\hline
\end{tabular}
\caption{Common SLIME and Sly commands}
\label{tab:slime}
\end{table}

While you type a call to a function, its argument list is shown at the
bottom of the window. Emacs indents Lisp code for you and shows which
parentheses match, so you read the structure of the code from its
indentation and never count parentheses.

\textbf{Try it:} Open \code{/tmp/hello.lisp} in Emacs, with SLIME
running. Put the cursor inside the definition of \code{hello} and
press \code{C-c C-c}. Switch to the REPL with \code{C-c C-z} and type
\code{(hello)}. Now go back to the file, change the string to
\code{"Hello, CL World!"}, press \code{C-c C-c} again, and call
\code{(hello)} again in the REPL. You have changed a running program
without restarting it.

\subsection{Other Editors}

Emacs is not required. Visual Studio Code has the Alive\index{Alive} extension,
which gives a REPL, evaluation from the editor, a debugger and an
inspector. Vim and Neovim have Vlime and Slimv. Lem\index{Lem}
(\url{https://lem-project.github.io}) is an Emacs-like editor written
in Common Lisp, which works with Lisp as soon as it starts. Allegro
CL and LispWorks each come with their own environment.

Whatever you choose, make sure it gives you the essentials: a REPL,
a key to send the current definition to the running Lisp, and
automatic indentation.

\section{The Init File}

When CL starts, it looks for an initialization file in your home
directory and loads it\index{init file}. In SBCL this is \code{.sbclrc}; in CCL it is
\code{.ccl-init.lisp}. The file holds ordinary Lisp, and is the place
to load the tools you want in every session. The next section puts
something in it.

\section{Libraries: Quicklisp and ASDF}

\subsection{Loading Libraries with Quicklisp}

Quicklisp\index{Quicklisp} (\url{https://www.quicklisp.org}) downloads and loads
open-source libraries, along with whatever they depend on. To install
it, download one file and load it:

\begin{verbatim}
curl -O https://beta.quicklisp.org/quicklisp.lisp
sbcl --load quicklisp.lisp
\end{verbatim}

\noindent
and then, at the Lisp prompt:

\begin{verbatim}
* (quicklisp-quickstart:install)
* (ql:add-to-init-file)
\end{verbatim}

The second form adds a few lines to your init file so that Quicklisp
is there whenever you start Lisp. From then on, one call fetches a
library and loads it:

\begin{verbatim}
CL-USER> (ql:quickload "alexandria")
To load "alexandria":
  Load 1 ASDF system:
    alexandria
; Loading "alexandria"

("alexandria")
CL-USER> (alexandria:flatten '(1 (2 (3 4)) 5))
(1 2 3 4 5)
\end{verbatim}

Alexandria, loaded here, is a collection of small utilities that most
other libraries use. \code{(ql:system-apropos "json")} searches the
library list by name. Chapter~4 uses several libraries, and
Chapter~5 says where to look for more.

Quicklisp is the common starting point, and there are other tools for
the same job, such as ocicl and Qlot, which add things like pinning
the exact versions a project uses.

\subsection{Defining Your Own System with ASDF}
\label{sec:asdf}

A program of any size is spread over several files, some of which have
to be loaded before others. ASDF\index{ASDF} is the standard tool for describing
that. It comes with every implementation mentioned in this chapter,
and it is what Quicklisp uses underneath. In ASDF, a library or
program is called a \emph{system}\index{system}, and is described by a file
ending in \code{.asd}.

Here is a complete small system, in a directory named
\code{greeter}. First \code{greeter.asd}:

\begin{verbatim}
(defsystem "greeter"
  :description "A very small example system."
  :version "0.1.0"
  :depends-on ("alexandria")
  :components ((:file "package")
               (:file "greet" :depends-on ("package"))))
\end{verbatim}

\noindent
then \code{package.lisp}:

\begin{verbatim}
(defpackage :greeter
  (:use :common-lisp)
  (:export #:greet))
\end{verbatim}

\noindent
and \code{greet.lisp}:

\begin{verbatim}
(in-package :greeter)

(defun greet (name)
  (format nil "Hello, ~a!" (string-capitalize name)))
\end{verbatim}

Put the \code{greeter} directory inside \code{\textasciitilde/common-lisp/},
which is where ASDF looks by default, and you can load the system by
name like any other library:

\begin{verbatim}
CL-USER> (ql:quickload "greeter")
To load "greeter":
  Load 1 ASDF system:
    greeter
; Loading "greeter"

("greeter")
CL-USER> (greeter:greet "reader")
"Hello, Reader!"
\end{verbatim}

ASDF compiles each file the first time and whenever it has changed,
keeps the fasl files out of your source directory, and loads the files
in an order that respects the dependencies you declared. Quicklisp
fetches Alexandria if it is not yet installed.

Start every project this way, however small. It costs three short
files, and it means anyone, including you a year from now, can load
the program with one line.

\section{Scripts and Standalone Programs}

\subsection{Running a File as a Script}

CL is normally used interactively, but it will also run a file from
the command line and exit\index{scripting}. Here is \code{script.lisp}, which copies
its input to its output with leading spaces removed from each line:

\begin{verbatim}
;;; Print each line of standard input with leading spaces removed.
(loop for line = (read-line *standard-input* nil)
      while line
      do (write-line (string-left-trim " " line)))
\end{verbatim}

\begin{verbatim}
sbcl --script script.lisp < datafile1 > datafile2
\end{verbatim}

A semicolon starts a comment\index{comments}, which runs to the end of the line. By
convention one semicolon is used for a comment after code on the same
line, two for a comment on its own line inside a definition, and three
for a comment at the top level of a file.

With \code{--script}, SBCL skips your init file, so Quicklisp is not
loaded and neither is ASDF. A script that needs them starts by
loading them itself.

\subsection{Building an Executable}

Because a running Lisp is an image\index{image} in memory, most implementations
can write that image to disk, and can write it as a program that
starts in a function of your choosing. This is how a Lisp application
is delivered to someone who does not have Lisp installed. In SBCL:

\begin{verbatim}
(defun main ()
  (format t "Hello from a standalone program!~%"))

(sb-ext:save-lisp-and-die "hello-app"
                          :toplevel #'main
                          :executable t)
\end{verbatim}

This writes a file \code{hello-app} that runs on any machine of the
same kind. It is about forty megabytes, because it contains the whole
Lisp, compiler included. Each implementation has its own function
for this. If you need one build recipe for several implementations,
ASDF can do it; look up \code{program-op} in its manual.

\section{Debugging in CL}
\label{sec:debugging}

Common Lisp has one of the best debugging and error-handling
arrangements of any language, and it is always on. There is no
separate debug mode and no debug build.

\subsection{The Debugger}

When an error happens in most languages, the program unwinds to the
top and dies, and you read a stack trace to work out what the state
must have been. When an error happens in Common Lisp, the program
\emph{stops where the error happened}, with everything still in place,
and you are put into the debugger\index{debugger}:

\begin{verbatim}
* (defun oops (x) (/ 10 x))
OOPS
* (oops 0)

debugger invoked on a DIVISION-BY-ZERO in thread
#<THREAD tid=59 "main thread" RUNNING {1200038003}>:
  arithmetic error DIVISION-BY-ZERO signalled
Operation was (/ 10 0).

Type HELP for debugger help, or (SB-EXT:EXIT) to exit from SBCL.

restarts (invokable by number or by possibly-abbreviated name):
  0: [ABORT] Exit debugger, returning to top level.

(SB-KERNEL::INTEGER-/-INTEGER 10 0)
0]
\end{verbatim}

The prompt \code{0]} is the debugger's. It is still a Lisp prompt:
you can evaluate any expression there, look at variables, and even
redefine the function that failed. \code{backtrace} shows the chain
of calls that led here. Typing \code{0}, or \code{abort}, takes the
restart numbered 0 and returns you to the normal prompt.

In SLIME or Sly the same information appears in its own window. Press
the number of a restart to choose it, \code{q} to abort back to the
REPL, Return on a line of the backtrace to see that call's local
variables, and \code{v} to jump to its source.

\subsection{Restarts}

The list headed ``restarts'' is the unusual part\index{restarts}. A restart is a
way of carrying on that the program itself has offered. Try opening
a file that does not exist:

\begin{verbatim}
* (with-open-file (in "/tmp/nope.conf") (read in))

debugger invoked on a FILE-DOES-NOT-EXIST in thread
#<THREAD tid=68 "main thread" RUNNING {1200038003}>:
  The file #P"/tmp/nope.conf" does not exist: No such file or directory

restarts (invokable by number or by possibly-abbreviated name):
  0: [CREATE   ] Reopen with :if-does-not-exist :create
  1: [CONTINUE ] Retry opening.
  2: [USE-VALUE] Try opening a different file.
  3: [ABORT    ] Exit debugger, returning to top level.
\end{verbatim}

You could now create the file from another window and choose restart
1, and the program would carry on from the middle as if nothing had
gone wrong. When the failure comes two hours into a long computation,
that is worth a great deal. Section~\ref{sec:conditions} shows how a
program handles errors itself, without the debugger.

\subsection{Stopping a Program Yourself}

To interrupt a computation that is taking too long, press Control-C in
the terminal, or \code{C-c C-c} in the SLIME REPL. You land in the
debugger, where you can look around and then either continue or abort.

To make a program stop at a particular place, call \code{break}\index{break@\texttt{break}} there:

\begin{verbatim}
(defun risky (n)
  (break "n is ~a" n)
  (* n 2))
\end{verbatim}

\noindent
Calling \code{(risky 21)} enters the debugger with the message ``n is
21'' and a \code{CONTINUE} restart, which resumes the function and
returns 42.

\subsection{Trace, Describe and Friends}

\code{Trace}\index{trace@\texttt{trace}} prints every call to a function and what it returned,
without your editing the function. It is especially good for seeing
how a recursive function works:

\begin{verbatim}
CL-USER> (defun our-length (list)
           (if (null list)
               0
               (+ (our-length (rest list)) 1)))
OUR-LENGTH
CL-USER> (trace our-length)
(OUR-LENGTH)
CL-USER> (our-length '(a b))
  0: (OUR-LENGTH (A B))
    1: (OUR-LENGTH (B))
      2: (OUR-LENGTH NIL)
      2: OUR-LENGTH returned 0
    1: OUR-LENGTH returned 1
  0: OUR-LENGTH returned 2
2
CL-USER> (untrace our-length)
T
\end{verbatim}

The running Lisp can also tell you about itself.
\code{Describe}\index{describe@\texttt{describe}} prints what is known about any object,
\code{documentation}\index{documentation@\texttt{documentation}} returns a documentation string, and
\code{apropos}\index{apropos@\texttt{apropos}} finds names containing a piece of text:

\begin{verbatim}
CL-USER> (documentation 'mapcar 'function)
"Apply FUNCTION to successive tuples of elements of LIST and MORE-LISTS.
Return list of FUNCTION return values."
CL-USER> (apropos "HASH-TABLE-" :cl)
HASH-TABLE-COUNT (fbound)
HASH-TABLE-P (fbound)
HASH-TABLE-REHASH-SIZE (fbound)
HASH-TABLE-REHASH-THRESHOLD (fbound)
HASH-TABLE-SIZE (fbound)
HASH-TABLE-TEST (fbound)
WITH-HASH-TABLE-ITERATOR (fbound)
\end{verbatim}

\subsection{Timing}

Wrap any expression in the \code{time}\index{time@\texttt{time}} macro to see how long it takes
and how much memory it allocates. The expression's value is returned
as usual:

\begin{verbatim}
CL-USER> (defun fib (n)
           (if (< n 2)
               n
               (+ (fib (- n 1)) (fib (- n 2)))))
FIB
CL-USER> (time (fib 30))
Evaluation took:
  0.023 seconds of real time
  0.021844 seconds of total run time (0.021844 user, 0.000000 system)
  95.65% CPU
  65,398,250 processor cycles
  0 bytes consed

832040
\end{verbatim}

``Consed'' means allocated; Section~\ref{sec:cons} explains the word.
For finding out where a larger program spends its time, the
implementations provide profilers. SBCL's is called
\code{sb-sprof}.

\section{Working with an AI Agent}
\label{sec:agents}

Many programmers now work with an AI coding agent\index{AI agents}: a language model
that reads the project, proposes changes, and runs commands. Lisp
suits this way of working unusually well, for the same reasons it
suits a person at a REPL.

An agent working in most languages writes code, runs the whole program
or its tests, and reads the output. An agent connected to a running
Lisp image can do what you do: evaluate one expression and look at the
answer, ask the image for a function's argument list or documentation
instead of guessing, redefine one function and try it, and inspect the
state left behind by an error. Mistakes are caught an expression at a
time. The connection is usually made with the Model Context
Protocol\index{Model Context Protocol} (MCP), a common way of giving a model tools to call. One
open-source bridge between MCP and a running Lisp is lisply-mcp,
\url{https://github.com/genworks/lisply-mcp}.

Some advice that has held up:

\begin{itemize}
\item Have the agent evaluate what it writes, in the image, before it
shows it to you. Code that has only been written is a guess.

\item Read what it produces. To check an agent's Lisp you need what
this book teaches: what gets evaluated and when, what a function
returns, what is modified in place.

\item Let it work in a Lisp that is safe to break, such as a
container, and keep your own work in version control.
\end{itemize}

An agent makes it more useful to understand the language, not less.
You are the one who decides whether the program is right.
